\section{Remove the analytic cutoff in the surplus ordinary-tail theorem}
The registered ordinary-tail theorem used u>=max(h,16).
The lower bound sixteen was used ONLY to invoke the analytic part
of the old ordinary quotient theorem. It is not needed in any of the
new cumulative, source-startup, bounded-cut or boundary arguments.

The registered degree-boundary theorem explicitly proves positive
GLOBAL ordinary LC of J at EVERY u=h>=2, including its independently
verified175 small-u contexts. Its sole [4],u3 exception concerns the
positive primitive of J, not ordinary LC of J.
For u>=h+1 use the registered full surplus quotient theorem.
Our ordinary coefficient argument needs only ordinary LC of J:
the positive primitive of the cumulative quotient B is J itself.

Every scalar startup/head guard in SurplusOrdinaryTail.lean requires
only w=u+r>=6. The source and negative-index/first-mode arguments have
no u>=16 dependency. Thus the identical written proof establishes
\[
 u\ge h,\quad u+r\ge6
 \quad\Longrightarrow\quad
 F_N^2\ge F_{N-1}F_{N+1}\quad(N\ge u+r+1)
\]
for EVERY finite padding and nonbinary capacity list.
The old175 contexts are REUSED, not recomputed/recounted or replaced
by new sampled hosts. There is NO new finite entrance premise.
In particular this applies to EVERY nonempty surplus prefix at u>=5.
The small w<=5 ordinary families are not claimed by this refinement;
the universal late-C proof uses the old small-binary theorem instead.
Relative ordinary LC and the Cassini-amplitude strengthening are still
separate OPEN assertions and are not used by the new first-crossing proof.

